\documentclass[12pt,a4paper]{article}

% Packages
\usepackage[utf8]{inputenc}
\usepackage[T1]{fontenc}
\usepackage[english]{babel}
\usepackage{mathtools}
\usepackage{amsmath}
\usepackage{amssymb}
\usepackage{amsthm}
\usepackage{geometry}
\usepackage{hyperref}
\usepackage{cleveref}
\usepackage{booktabs}
\usepackage{array}
\usepackage{xcolor}
\usepackage{float}

\geometry{a4paper, margin=2.5cm}

\hypersetup{
  colorlinks=true,
  linkcolor=blue!70!black,
  citecolor=green!60!black,
  urlcolor=blue!60!black,
  pdftitle={Mathematical Analysis of the Cone Surface},
  pdfauthor={Dogan Balban},
  pdfkeywords={cone, surface of revolution, Gaussian curvature, surface area, volume}
}

% Theorem environments
\theoremstyle{plain}
\newtheorem{theorem}{Theorem}[section]
\newtheorem{lemma}[theorem]{Lemma}
\newtheorem{proposition}[theorem]{Proposition}
\newtheorem{corollary}[theorem]{Corollary}

\theoremstyle{definition}
\newtheorem{definition}{Definition}[section]
\newtheorem{example}{Example}[section]

\theoremstyle{remark}
\newtheorem{remark}{Remark}[section]

% cleveref names
\crefname{theorem}{Theorem}{Theorems}
\crefname{proposition}{Proposition}{Propositions}
\crefname{corollary}{Corollary}{Corollaries}
\crefname{definition}{Definition}{Definitions}
\crefname{remark}{Remark}{Remarks}
\crefname{equation}{Eq.}{Eqs.}
\crefname{table}{Table}{Tables}
\crefname{section}{Section}{Sections}

\newcommand{\R}{\mathbb{R}}
\newcommand{\vek}[1]{\vec{#1}}

\title{\textbf{Mathematical Analysis of the Cone Surface}\\[0.5em]
\large Complete Differential-Geometric Treatment:\\
Curvature, Surface Area, Volume, and Geometric Properties}
\author{Dogan Balban}
\date{\today}

\begin{document}
\maketitle

\begin{abstract}
This paper presents a complete mathematical analysis of the cone surface defined by
\[
z = -\frac{1}{12}\cdot 3^{2} + 6\cdot\sqrt{\frac{x^{2}}{\bigl(\tfrac{12}{\pi}\bigr)^{2}}
+\frac{y^{2}}{\bigl(\tfrac{12}{\pi}\bigr)^{2}}}.
\]
We investigate the algebraic structure, geometric properties, differential geometry,
curvature behaviour, and surface parametrization. The analysis covers partial derivatives,
gradient, normal vectors, surface area, volume, and a complete characterisation of this
cone.
\end{abstract}

\tableofcontents
\newpage

%=============================================================
\section{Introduction and Simplification}

\subsection{Given Expression}

We consider the surface defined by:
\begin{equation}\label{eq:original}
z = -\frac{1}{12}\cdot 3^{2} + 6\cdot\sqrt{\frac{x^{2}}{\bigl(\tfrac{12}{\pi}\bigr)^{2}}
+\frac{y^{2}}{\bigl(\tfrac{12}{\pi}\bigr)^{2}}}
\end{equation}

\subsection{Algebraic Simplification}

\begin{proposition}[Simplification]\label{prop:simplification}
Expression \eqref{eq:original} simplifies to:
\begin{equation}\label{eq:simplified}
z = -\frac{3}{4} + \frac{\pi}{2}\sqrt{x^2 + y^2}.
\end{equation}
\end{proposition}

\begin{proof}
\textbf{Step~1} (Constant term): $-\tfrac{1}{12}\cdot 3^{2} = -\tfrac{9}{12} = -\tfrac{3}{4}$.

\textbf{Step~2} (Radicand): $\dfrac{x^{2}}{(12/\pi)^{2}} = x^{2}\cdot\dfrac{\pi^2}{144}$,
and analogously for $y$.

\textbf{Step~3}: $\sqrt{\dfrac{\pi^2}{144}(x^2+y^2)} = \dfrac{\pi}{12}\sqrt{x^2+y^2}$.

\textbf{Step~4}: $6\cdot\dfrac{\pi}{12}\sqrt{x^2+y^2} = \dfrac{\pi}{2}\sqrt{x^2+y^2}$.

\textbf{Step~5}: Combining: $z = -\tfrac{3}{4}+\tfrac{\pi}{2}\sqrt{x^2+y^2}$.
\end{proof}

\subsection{Cylindrical Coordinates}

\begin{definition}[Cylindrical form]\label{def:cylindrical}
Setting $x = r\cos\theta$, $y = r\sin\theta$, $r = \sqrt{x^2+y^2}$, the equation
reduces to:
\begin{equation}\label{eq:cylindrical}
\boxed{z = -\frac{3}{4} + \frac{\pi}{2}\,r.}
\end{equation}
\end{definition}

\begin{remark}
Equation~\eqref{eq:cylindrical} is independent of $\theta$: the surface is a
\textbf{surface of revolution} about the $z$-axis -- specifically, a \textbf{right
circular cone}.
\end{remark}

%=============================================================
\section{Geometric Characterisation}

\subsection{Identification as a Cone}

\begin{theorem}[Cone property]\label{thm:cone}
The surface defined by~\eqref{eq:cylindrical} is a right circular cone with:
\begin{itemize}
\item apex at $(0,0,-\tfrac{3}{4})$,
\item axis along the positive $z$-axis,
\item slope $m = \tfrac{\pi}{2}$.
\end{itemize}
\end{theorem}

\begin{proof}
The equation $z = -\tfrac{3}{4}+\tfrac{\pi}{2}r$ is linear in $r$ with positive
coefficient. At $r=0$ we have $z=-\tfrac{3}{4}$ (apex). For $r>0$, $z$ increases
linearly -- the characteristic property of a right circular cone.
\end{proof}

\subsection{Cone Angles}\label{sec:cone_angle}

\begin{definition}[Cone half-angle]\label{def:cone_angle}
The \emph{cone half-angle} $\alpha$ is the angle between the cone axis ($z$-axis)
and a generator (slant line).
\end{definition}

\begin{proposition}[Cone half-angle]\label{prop:cone_angle}
For our cone:
\begin{equation}\label{eq:cone_angle}
\alpha = \arctan\!\left(\frac{2}{\pi}\right) \approx 0.5669\text{ rad} \approx 32.48^\circ.
\end{equation}
\end{proposition}

\begin{proof}
A generator runs from $(r_0,\theta_0,z_0)$ to $(r_0+\Delta r,\theta_0,z_0+\Delta z)$
with $\Delta z = \tfrac{\pi}{2}\,\Delta r$. The angle \emph{between the axis and the
generator} satisfies:
\[
\tan\alpha = \frac{\Delta r}{\Delta z} = \frac{1}{\pi/2} = \frac{2}{\pi}
\quad\Rightarrow\quad
\alpha = \arctan\!\left(\frac{2}{\pi}\right) \approx 32.48^\circ.
\]
\end{proof}

\begin{remark}[Elevation angle and complementary angle]\label{rem:angles}
The complementary \emph{elevation angle} (angle between generator and the horizontal
$xy$-plane) is:
\[
\beta = \arctan\!\left(\frac{\pi}{2}\right) \approx 1.0039\text{ rad} \approx 57.52^\circ,
\]
with $\alpha+\beta = \tfrac{\pi}{2}$ ($32.48^\circ+57.52^\circ = 90^\circ$).
Some literature uses $\beta$ as the ``cone angle''; in this paper $\alpha$ always
denotes the angle between axis and generator (\cref{def:cone_angle}).
\end{remark}

\begin{corollary}[Apex angle]\label{cor:apex_angle}
The full apex angle (between two opposite generators) is:
\[
2\alpha = 2\arctan\!\left(\frac{2}{\pi}\right) \approx 64.96^\circ.
\]
\end{corollary}

\subsection{Cross-Sections}

\begin{proposition}[Base circle]\label{prop:base_circle}
The cone intersects $z=0$ in a circle of radius
$r_0 = \tfrac{3}{2\pi}\approx 0.4775$.
\end{proposition}

\begin{proposition}[Horizontal cross-section]\label{prop:horizontal}
The intersection with $z=h$ (for $h\geq -\tfrac{3}{4}$) is a circle of radius:
\begin{equation}\label{eq:radius_h}
r(h) = \frac{2}{\pi}\!\left(h+\frac{3}{4}\right).
\end{equation}
\end{proposition}

%=============================================================
\section{Functional Representation}

\subsection{Explicit Form and Partial Derivatives}

As a graph: $f(x,y) = -\tfrac{3}{4}+\tfrac{\pi}{2}\sqrt{x^2+y^2}$, $(x,y)\in\R^2$.

\begin{proposition}[Partial derivatives]\label{prop:partial}
For $(x,y)\neq(0,0)$:
\begin{align}
f_x &= \frac{\pi x}{2r},\quad f_y = \frac{\pi y}{2r};\label{eq:grad1}\\
f_{xx} &= \frac{\pi y^2}{2r^3},\quad f_{yy} = \frac{\pi x^2}{2r^3},\quad
f_{xy} = -\frac{\pi xy}{2r^3}.\label{eq:grad2}
\end{align}
\end{proposition}

\begin{remark}
At the apex $(0,0)$ the partial derivatives are undefined: this is the geometric
singularity of the cone.
\end{remark}

\subsection{Gradient and Normal Vector}

\begin{proposition}[Gradient and normal]\label{prop:gradient}
\begin{align*}
\nabla f &= \left(\frac{\pi x}{2r},\;\frac{\pi y}{2r}\right),\quad
|\nabla f| = \frac{\pi}{2}\quad(\text{constant for all }(x,y)\neq(0,0)),\\
\vek{n} &= \left(\frac{\pi x}{2r},\;\frac{\pi y}{2r},\;-1\right),\quad
\hat{n} = \frac{1}{\sqrt{1+\pi^2/4}}\vek{n}.
\end{align*}
\end{proposition}

%=============================================================
\section{Parametrization}

\begin{definition}[Parametrization]\label{def:param}
A natural parametrization of the cone surface is:
\begin{equation}\label{eq:parametrization}
\vek{\gamma}(r,\theta) = \bigl(r\cos\theta,\;r\sin\theta,\;{-\tfrac{3}{4}+\tfrac{\pi}{2}r}\bigr)^\top,
\quad r\geq 0,\;\theta\in[0,2\pi).
\end{equation}
\end{definition}

\begin{proposition}[Tangent vectors and cross product]\label{prop:cross}
\begin{align*}
\frac{\partial\vek{\gamma}}{\partial r} &= (\cos\theta,\;\sin\theta,\;\tfrac{\pi}{2})^\top,\quad
\frac{\partial\vek{\gamma}}{\partial\theta} = (-r\sin\theta,\;r\cos\theta,\;0)^\top,\\
\frac{\partial\vek{\gamma}}{\partial r}\times\frac{\partial\vek{\gamma}}{\partial\theta}
&= \bigl(-\tfrac{\pi r}{2}\cos\theta,\;-\tfrac{\pi r}{2}\sin\theta,\;r\bigr)^\top,\\
\left|\frac{\partial\vek{\gamma}}{\partial r}\times\frac{\partial\vek{\gamma}}{\partial\theta}\right|
&= r\sqrt{1+\frac{\pi^2}{4}}.
\end{align*}
\end{proposition}

\begin{proof}
Direct computation. Magnitude:
$\sqrt{(\pi r/2)^2(\cos^2\theta+\sin^2\theta)+r^2} = r\sqrt{1+\pi^2/4}$.
\end{proof}

%=============================================================
\section{Surface Area}

\begin{theorem}[Lateral surface area of frustum]\label{thm:surface_area}
The lateral surface area of the cone frustum between radii $r_1$ and $r_2$
($0\leq r_1<r_2$) is:
\begin{equation}\label{eq:surface_area}
A = \pi\sqrt{1+\frac{\pi^2}{4}}\cdot(r_2^2 - r_1^2).
\end{equation}
\end{theorem}

\begin{proof}
$A = \int_0^{2\pi}\!\int_{r_1}^{r_2}\sqrt{1+\pi^2/4}\cdot r\,dr\,d\theta
= \sqrt{1+\pi^2/4}\cdot 2\pi\cdot\tfrac{r_2^2-r_1^2}{2}
= \pi\sqrt{1+\pi^2/4}\cdot(r_2^2-r_1^2)$.
\end{proof}

\begin{corollary}[Surface area from apex]\label{cor:surface_from_apex}
The lateral surface area from the apex to radius $r$:
\begin{equation}\label{eq:surface_apex}
A(r) = \pi r^2\sqrt{1+\frac{\pi^2}{4}} \approx 1.8621\cdot\pi r^2.
\end{equation}
\end{corollary}

%=============================================================
\section{Volume}

\begin{theorem}[Volume of frustum]\label{thm:volume}
The volume of the frustum between heights $h_1$ and $h_2$
($-\tfrac{3}{4}\leq h_1<h_2$) is:
\begin{equation}\label{eq:volume}
V = \frac{4}{3\pi}\left[\left(h_2+\frac{3}{4}\right)^3 - \left(h_1+\frac{3}{4}\right)^3\right].
\end{equation}
\end{theorem}

\begin{proof}
Using $r(h) = \tfrac{2}{\pi}(h+\tfrac{3}{4})$ from \cref{prop:horizontal}:
\[
V = \int_{h_1}^{h_2}\pi\,[r(h)]^2\,dh
= \frac{4}{\pi}\int_{h_1}^{h_2}\!\left(h+\frac{3}{4}\right)^2 dh
= \frac{4}{3\pi}\!\left[\left(h+\frac{3}{4}\right)^3\right]_{h_1}^{h_2}.
\]
\textit{Verification}: For $h_1=-\tfrac{3}{4}$, $h_2=h$, the classical cone formula
$V=\tfrac{1}{3}\pi R^2 H$ with $R=r(h)$, $H=h+\tfrac{3}{4}$ gives the same result.
\end{proof}

%=============================================================
\section{Curvature Analysis}

\subsection{Gaussian Curvature}

\begin{theorem}[Gaussian curvature]\label{thm:gauss}
At every regular point of the cone:
\[
K = 0.
\]
\end{theorem}

\begin{proof}
A cone is a \emph{developable surface} (locally isometric to the plane): $K=0$.
Directly: $\kappa_1 = 0$ along generators (straight lines), so $K = \kappa_1\kappa_2 = 0$.
\end{proof}

\subsection{Mean Curvature}

\begin{proposition}[Mean curvature]\label{prop:mean_curvature}
At a point with radial coordinate $r>0$:
\begin{equation}\label{eq:mean_curvature}
H = \frac{\pi}{4r\sqrt{1+\dfrac{\pi^2}{4}}} = \frac{\pi}{2r\sqrt{4+\pi^2}}.
\end{equation}
\end{proposition}

\begin{proof}
For $z = f(x,y)$, the mean curvature is:
\[
H = \frac{(1+f_y^2)f_{xx} - 2f_xf_yf_{xy} + (1+f_x^2)f_{yy}}{2(1+f_x^2+f_y^2)^{3/2}}.
\]
Since $f_x^2+f_y^2 = \pi^2/4$ (constant, \cref{prop:gradient}), the denominator
is $2(1+\pi^2/4)^{3/2}$. Expanding the numerator using \cref{prop:partial}:
\[
\text{Numerator} = \frac{\pi y^2}{2r^3}+\frac{\pi x^2}{2r^3}
+\frac{\pi^3(x^4+y^4+2x^2y^2)}{8r^5}
= \frac{\pi}{2r}+\frac{\pi^3}{8r} = \frac{\pi(4+\pi^2)}{8r}.
\]
Hence:
\[
H = \frac{\pi(4+\pi^2)/8r}{2[(4+\pi^2)/4]^{3/2}}
= \frac{\pi}{2r\sqrt{4+\pi^2}} = \frac{\pi}{4r\sqrt{1+\pi^2/4}}. \qed
\]
\end{proof}

\begin{remark}
Equivalently, using the parametric formula $H = m/(2r\sqrt{1+m^2})$ with slope $m=\pi/2$
yields $H = (\pi/2)/(2r\sqrt{1+\pi^2/4}) = \pi/(4r\sqrt{1+\pi^2/4})$ directly.
Since $\kappa_1=0$, the second principal curvature is $\kappa_2 = 2H$.
\end{remark}

%=============================================================
\section{Geodesics}

\begin{theorem}[Geodesics on the cone]
The geodesic lines on the cone are:
\begin{enumerate}
\item \textbf{Generators}: straight lines radially from the apex.
\item \textbf{Rolled geodesics}: curves that become straight lines when the cone
is developed (unrolled) into a planar sector.
\end{enumerate}
The developed cone is a circular sector with opening angle
$2\pi\sin\alpha \approx 2\pi\times 0.5365 \approx 3.370\text{ rad} \approx 193.1^\circ$.
\end{theorem}

%=============================================================
\section{Numerical Examples}

\begin{table}[H]
\centering
\caption{Cone parameters}\label{tab:parameters}
\begin{tabular}{lcc}
\toprule
Parameter & Exact value & Numerical \\
\midrule
Apex ($z$-coordinate) & $-\tfrac{3}{4}$ & $-0.750$ \\
Slope $m$ & $\tfrac{\pi}{2}$ & $1.5708$ \\
Half-angle $\alpha$ (axis--generator) & $\arctan(2/\pi)$ & $0.5669\text{ rad} = 32.48^\circ$ \\
Elevation angle $\beta$ (generator--horiz.) & $\arctan(\pi/2)$ & $1.0039\text{ rad} = 57.52^\circ$ \\
Apex angle $2\alpha$ & $2\arctan(2/\pi)$ & $64.96^\circ$ \\
Base radius ($z=0$) & $\tfrac{3}{2\pi}$ & $0.4775$ \\
$\sqrt{1+\pi^2/4}$ & $\tfrac{\sqrt{4+\pi^2}}{2}$ & $1.8621$ \\
\bottomrule
\end{tabular}
\end{table}

\begin{table}[H]
\centering
\caption{Lateral surface area: $A = \pi\sqrt{1+\pi^2/4}\cdot(r_2^2-r_1^2)$}
\label{tab:surface_area}
\begin{tabular}{ccc}
\toprule
$r_1$ & $r_2$ & $A$ \\
\midrule
$0$ & $5$  & $146.249$ \\
$0$ & $10$ & $584.995$ \\
$5$ & $10$ & $438.746$ \\
$0$ & $20$ & $2339.979$ \\
\bottomrule
\end{tabular}
\end{table}

\begin{table}[H]
\centering
\caption{Volume: $V = \tfrac{4}{3\pi}\bigl[(h_2+\tfrac{3}{4})^3-(h_1+\tfrac{3}{4})^3\bigr]$}
\label{tab:volume}
\begin{tabular}{ccc}
\toprule
$h_1$ & $h_2$ & $V$ \\
\midrule
$-0.75$ & $0$  & $0.179$ \\
$0$     & $10$ & $527.068$ \\
$-0.75$ & $10$ & $527.247$ \\
$0$     & $20$ & $3791.601$ \\
\bottomrule
\end{tabular}
\end{table}

%=============================================================
\section{Summary of All Properties}

\begin{table}[H]
\centering
\caption{Complete summary of cone properties}\label{tab:summary}
\begin{tabular}{ll}
\toprule
Property & Value / Formula \\
\midrule
Simplified form & $z = -\tfrac{3}{4}+\tfrac{\pi}{2}r$ \\
Surface type & Right circular cone \\
Apex & $(0,0,-\tfrac{3}{4})$ \\
Slope & $m = \tfrac{\pi}{2}$ \\
Half-angle $\alpha$ (axis--generator) & $\arctan(2/\pi)\approx 32.48^\circ$ \\
Elevation angle $\beta$ & $\arctan(\pi/2)\approx 57.52^\circ$ \\
Apex angle $2\alpha$ & $2\arctan(2/\pi)\approx 64.96^\circ$ \\
Gaussian curvature & $K=0$ (developable) \\
Mean curvature & $H = \dfrac{\pi}{4r\sqrt{1+\pi^2/4}}$ \\[4pt]
Gradient magnitude & $|\nabla f| = \tfrac{\pi}{2}$ \\
Lateral surface area & $A = \pi\sqrt{1+\pi^2/4}\cdot(r_2^2-r_1^2)$ \\
Volume & $V = \tfrac{4}{3\pi}\bigl[(h_2+\tfrac{3}{4})^3-(h_1+\tfrac{3}{4})^3\bigr]$ \\
\bottomrule
\end{tabular}
\end{table}

%=============================================================
\section{Conclusion}

This paper has presented a complete mathematical analysis of the cone surface
$z = -\tfrac{3}{4}+\tfrac{\pi}{2}r$:

\begin{enumerate}
\item The complex original expression simplifies elegantly to $z = -\tfrac{3}{4}+\tfrac{\pi}{2}r$.
\item The surface is a right circular cone with apex $(0,0,-\tfrac{3}{4})$ and slope $\pi/2$.
\item Gaussian curvature $K=0$ (developable surface, locally isometric to the plane).
\item Cone half-angle (axis--generator): $\alpha = \arctan(2/\pi)\approx 32.48^\circ$;
elevation angle: $\beta = \arctan(\pi/2)\approx 57.52^\circ$; apex angle: $2\alpha\approx 64.96^\circ$.
\item Gradient is constant: $|\nabla f| = \pi/2$.
\item Mean curvature: $H = \pi/(4r\sqrt{1+\pi^2/4})$, decreasing as $r\to\infty$.
\item Closed-form expressions exist for lateral surface area and volume.
\end{enumerate}

%=============================================================

\begin{thebibliography}{9}

\bibitem{docarmo}
M.~P.~do~Carmo,
\textit{Differential Geometry of Curves and Surfaces}, revised ed.
Dover Publications, Mineola, NY, 2016.

\bibitem{struik}
D.~J.~Struik,
\textit{Lectures on Classical Differential Geometry}, 2nd~ed.
Dover Publications, New York, 1988.

\bibitem{gray}
A.~Gray, E.~Abbena, and S.~Salamon,
\textit{Modern Differential Geometry of Curves and Surfaces with Mathematica}, 3rd~ed.
Chapman \& Hall/CRC, Boca Raton, FL, 2006.

\bibitem{pressley}
A.~Pressley,
\textit{Elementary Differential Geometry}, 2nd~ed.
Springer, London, 2010.

\bibitem{balban_helix}
D.~Balban,
\textit{Differential-Geometric Analysis of the Conical Helix},
Independent mathematical publication, 2026.

\end{thebibliography}

\end{document}
